\section{Approximate zeta functions and the Euler product}\label{ki:sec:approx}

\subsection{Ki's approximation theorems}

Ki's 2005 paper truncates the Chowla Selberg expansion of an Epstein zeta function \cite{Ki05}. The partial sums keep the functional equation and have no Euler product, and Ki proves that in every strip containing the critical line all but finitely many of their zeros are simple and on the line. Later papers do the same for modified Epstein zeta functions \cite{Ki06a} and for approximations of the Ramanujan $\Xi$ function attached to $\Delta$ \cite{Ki08a}. % UNCHECKED: the statements of Ki05, Ki06a and Ki08a, and the sums of the Riemann zeta function of Ki, are quoted from the chapter source
Lemma~\ref{ki:lem:dict}(a) turns each of these into a Thorin statement once the approximation $F$ is even, real on the real axis, of order less than two in the variable, has $F(0)\ne0$, and all but finitely many of all its zeros, not only those in a fixed strip, are real; positivity of a kernel is not needed for this step, and where the approximation is not known to be a Fourier transform of a positive kernel Lemma~\ref{ki:lem:dict}(b) does not apply and these hypotheses must be checked. The conclusion is then that the reciprocal $F(0)/F(i\sqrt s)$ is eventually Thorin in the sense of Proposition~\ref{ki:prop:eventual}: a GGC transform up to a finite rational factor. The truncations and approximations $\zeta_N$ of Section~\ref{ch:rz4} and the friable approximant of Chapter~\ref{chap:miyatake} belong to the same family, and Ki's method, an argument principle count set against a count of real zeros by sign changes, is a template for proving eventual Thorin statements about them. % UNCHECKED: description of Ki's method

\subsection{The Davenport and Heilbronn obstruction}

Davenport and Heilbronn \cite{DH36} showed that the Epstein zeta function of a positive definite binary quadratic form of class number greater than one has zeros off the critical line (in the half plane $\re s>1$). Such a function has a functional equation of Riemann type, and its completed form has the inversion $t\mapsto1/t$ of a theta function. Its reciprocal Xi function is not a GGC, by Lemma~\ref{ki:lem:dict}(a).

\begin{remark}[What a Thorin proof must use, heuristic]\label{ki:rem:euler}
Any property shared by $\xi$ and the Davenport Heilbronn functions cannot imply RH. The functional equation and the inversion of the theta function are shared. If the P\'olya type kernel of such an Epstein zeta function is positive and has the same doubly exponential decay as $\Phi$, which is not established here, then positivity and decay of the kernel are shared as well. % UNCHECKED: positivity and decay of the Polya type kernel of the Epstein zeta function of a form of class number greater than one
An argument that derives a positive Thorin measure for $1/\Xi$ from these properties alone would then apply equally to those functions and would be invalid. This indicates that a correct argument must use the Euler product, or another property that separates $\zeta$ from them. It is a heuristic principle, not a theorem.
\end{remark}

The principle supports placing the Euler product L\'evy representation, which writes $\log(1/\zeta)$ through the primes (Section~\ref{bp:sec:levy}, Section~\ref{ch:arith}) and feeds that L\'evy measure to the sine squared lemma, at the centre of the programme. Proposition~\ref{ki:prop:obstruction} and Remark~\ref{ki:rem:euler} point the same way: the Gaussian tilt keeps the positivity, evenness and decay of the kernel but not the theta inversion, while the passage to Epstein functions keeps the functional equation and the theta inversion, and positivity and decay only under the unproved hypothesis of the remark. An argument blind to the Euler product cannot work.

\subsection{Weng zeta functions as a calibration set}

Weng's high rank zeta functions are built from shifted completed zeta values with a functional equation of Riemann type; they have no Euler product. Ki \cite{Ki10}, and Ki, Komori and Suzuki \cite{KKS15} for Chevalley groups, proved that their nontrivial zeros lie on the critical line. % UNCHECKED: scope of Ki10 and KKS15 (all ranks, all Chevalley groups, or all but finitely many zeros), and the construction of Weng's functions from completed zeta values
If that holds, the reciprocals of the corresponding Xi functions have positive Thorin measure by Lemma~\ref{ki:lem:dict}(a), and they are a test set for the numerical pipeline of this book (Chapter~\ref{chap:measure}), not for an argument that uses the Euler product: compute the Thorin measure through the sine squared lemma, confirm positivity and measure the margin. Because they have no Euler product, they are a different test from the one that Remark~\ref{ki:rem:euler} describes.

\section{Zeros of the derivative of zeta, open problems and notes}\label{ki:sec:open}

Ki showed, on RH, that $\liminf(\beta'-\tfrac12)\log\gamma'\ne0$ over the zeros $\beta'+i\gamma'$ of $\zeta'$ holds exactly when $\zeta'/\zeta(s)=1/(s-\rho)+O(\log t)$ near the line, with $\rho$ the nearest zero \cite{Ki08b}. Farmer and Ki tie zeros of $\zeta'$ near the line to small gaps between zeros of $\zeta$ and to Landau Siegel zeros \cite{FarmerKi12}, and Gonek and Ki study the pair correlation of the zeros of the real and imaginary parts of $\zeta$ \cite{GonekKi18}. % UNCHECKED: the statements of Ki08b, FarmerKi12 and GonekKi18 are quoted from the chapter source
In Thorin language $\zeta'/\zeta$ is the derivative of the L\'evy exponent, so Ki's criterion is a local regularity condition on the L\'evy measure near the critical line, and in the Farmer and Ki dichotomy small gaps appear as nearly coalescing Thorin atoms. These results concern the spacing of the zeros (Section~\ref{ch:spacing}); they give no new route to positivity.

For the programme the chapter yields the following, in order of use. Proposition~\ref{ki:prop:obstruction} (proved) says that no property of $\Phi$ that survives a Gaussian tilt can carry a proof of RH. Remark~\ref{ki:rem:euler} (heuristic) says that an argument blind to the Euler product cannot work, and Ki's approximation theorems and the Weng functions supply test cases for the numerical pipeline. Proposition~\ref{ki:prop:bm} (proved under RH, with Ki's theorem as input) and Remark~\ref{ki:conj:rate} (conjecture, computed) say that under RH the scaled clocks of the high derivatives converge to the Brownian exit time and conjecturally their atoms move, to first order, only through the third central moment. Proposition~\ref{ki:prop:eventual} and Corollary~\ref{ki:cor:parab} (proved, from the theorem of Ki, Kim and Lee and de Bruijn's strip theorem) put the theorem of Ki, Kim and Lee and the confocal parabola picture in Thorin form.

\begin{problem}[Prove the Edgeworth rate]\label{ki:prob:rate}
Prove the conjecture of Remark~\ref{ki:conj:rate}, with explicit error terms, by a saddle point analysis of $\ell_n$ and Lemma~\ref{ki:lem:zeroeq}. The local theorem of Ki and Kim (Section~\ref{pol:sec:papers}) gives that for every $R$ the zeros of $\Xi^{(n)}$ in $|\theta|<R$ are real for large $n$, so that the number $K_n$ of atoms of $\psi_n$ known to be real tends to infinity; determine how large $K_n$ may be unconditionally. The expansion is governed by the condition $c_{K_n}^3|\kappa_n|/m_n^3\to0$, that is $K_n=o\bigl((n^2\log n)^{1/3}\bigr)$; reality of the atoms is a separate question.
\end{problem}

\begin{problem}[Descending the derivative ladder]\label{ki:prob:ladder}
Find an explicit $n_0$ and heights $T(n)\to\infty$ such that all zeros of $\Xi^{(2n)}$ below $T(n)$ are real for $n\ge n_0$. This is the effective form of the local theorem of Ki and Kim \cite{KiKim00} and of Problem~\ref{pol:op:deriv}. It would give a sequence of certified positive Thorin measures computed from the moments of $\Phi$ alone, and a derivative analogue of the barrier.
\end{problem}

\begin{problem}[A property that breaks the obstruction]\label{ki:prob:property}
Identify a property of $\Phi$ used by the sine squared lemma that is destroyed by $\Phi\mapsto\Phi e^{-bu^2}$ for every $b>0$ and that fails for the P\'olya kernels of Epstein zeta functions of class number greater than one.
\end{problem}

\begin{problem}[Weng calibration]\label{ki:prob:weng}
Compute the Thorin measures of reciprocal Weng Xi functions of rank two and three through the sine squared lemma and compare their margin of positivity with that of $\sum_{\gamma>0}\delta_{\gamma^2}$.
\end{problem}

\begin{problem}[Eventual Thorin for the friable approximant]\label{ki:prob:friable}
Adapt Ki's argument principle and sign change method to prove that the reciprocal of the friable approximant of Chapter~\ref{chap:miyatake} is eventually Thorin, with an explicit bound on the exceptional set.
\end{problem}

\paragraph{Notes.} The theorems on the heat flow and on repeated differentiation are due to Ki and his coauthors as cited, to de Bruijn, Newman and to Rodgers and Tao; Proposition~\ref{ki:prop:obstruction} is the theorem of Rodgers and Tao read through the Thorin dictionary, in the form of Proposition~\ref{prop:biane:damping}(a), and Remark~\ref{ki:rem:euler} rests on the Davenport Heilbronn functions. The computations of Section~\ref{ki:sec:numbers} are those of \texttt{code/ki\_tilt.py} and \texttt{code/ki\_check.py} (data in \texttt{data/ki\_tilt\_book.json}, \texttt{data/ki\_tilt\_chapter.json} and \texttt{data/ki\_check.log}), which compute the tilted laws in mpmath and take about two seconds for each $n$.
